\documentclass[11pt]{article}

\usepackage[utf8]{inputenc}
\usepackage[T1]{fontenc}
\usepackage{amsmath,amssymb,amsthm}
\usepackage[margin=1.1in]{geometry}
\usepackage{hyperref}

\newtheorem{theorem}{Theorem}
\newtheorem{proposition}[theorem]{Proposition}
\newtheorem{lemma}[theorem]{Lemma}
\theoremstyle{remark}
\newtheorem{remark}[theorem]{Remark}

\title{Asymptotics of the Gram function and of its first derivative\\
  from the implicit definition $\vartheta(t_u)=(u-1)\pi$}
\author{}
\date{}

\begin{document}
\maketitle

\begin{abstract}
  We derive the classical asymptotic expansions for the Gram function
  $t_u$ (Lavrik) and its first derivative $t_u'$ (Korolev) --- equations
  (8) and (9) of \emph{Higher derivatives of the Gram function}
  (Dundulis, Garunk\v{s}tis, Laurin\v{c}ikas, \v{S}im\.{e}nas) --- directly
  from the implicit definition $\vartheta(t_u) = (u-1)\pi$ (equation (7)
  there), using only two standard facts about the Riemann--Siegel theta
  function $\vartheta$.
\end{abstract}

\section{Setup and input facts}

Throughout, $u \to +\infty$, and we abbreviate
\[
  x_u \;:=\; \frac{\log\log u}{\log u} \;\longrightarrow\; 0 .
\]

The Riemann--Siegel theta function $\vartheta(t)$ is the continuous branch
of $\arg\bigl(\pi^{-s/2}\Gamma(s/2)\bigr)$ along the segment from
$s = 1/2$ to $s = 1/2 + it$. We use the following two facts, both of
which follow from the Karatsuba--Korolev decomposition
$\vartheta(t) = \delta(t) - \pi\varphi(t) - \pi$ with
$\varphi(t) = -\tfrac{t}{2\pi}\log\tfrac{t}{2\pi} + \tfrac{t}{2\pi} -
\tfrac{7}{8}$ and $\delta^{(n)}(t) = O(t^{-n-1})$ for $n \ge 0$:
\begin{align}
  \vartheta(t)  &= \frac{t}{2}\,\log\frac{t}{2\pi} - \frac{t}{2}
                  - \frac{\pi}{8} + O\!\left(\frac{1}{t}\right),
                  & t \to +\infty, \tag{$\vartheta$1}\label{eq:theta1}\\[2pt]
  \vartheta'(t) &= \frac{1}{2}\,\log\frac{t}{2\pi}
                  + O\!\left(\frac{1}{t^{2}}\right),
                  & t \to +\infty. \tag{$\vartheta$2}\label{eq:theta2}
\end{align}
By \eqref{eq:theta2}, $\vartheta'(t) > 0$ for $t \ge 7$, so $\vartheta$
is strictly increasing on $[7,\infty)$ and tends to $+\infty$ there by
\eqref{eq:theta1}. Hence, for $u \ge \vartheta(7)/\pi + \pi$, the
\emph{Gram function} $t_u \ge 7$ is well defined by the implicit
equation
\begin{equation}
  \vartheta(t_u) \;=\; (u-1)\,\pi, \tag{7}\label{eq:seven}
\end{equation}
$u \mapsto t_u$ is increasing, and $t_u \to +\infty$ as $u \to +\infty$.

Our goal is to prove the following two expansions.

\begin{proposition}[Lavrik; equation (8)]\label{prop:eight}
  As $u \to +\infty$,
  \begin{equation}
    t_u \;=\; \frac{2\pi u}{\log u}
      \left(1 + (1+o(1))\,\frac{\log\log u}{\log u}\right).
    \tag{8}\label{eq:eight}
  \end{equation}
\end{proposition}

\begin{proposition}[Korolev; equation (9)]\label{prop:nine}
  As $u \to +\infty$,
  \begin{equation}
    \frac{dt_u}{du} \;=\; \frac{2\pi}{\log u}
      \left(1 + (1+o(1))\,\frac{\log\log u}{\log u}\right).
    \tag{9}\label{eq:nine}
  \end{equation}
\end{proposition}

\section{An algebraic form of the implicit equation}

Substituting \eqref{eq:theta1} into \eqref{eq:seven} and multiplying by
$2$ gives
\[
  t_u \log\frac{t_u}{2\pi} - t_u
  \;=\; 2\pi u - 2\pi + \frac{\pi}{4} + O\!\left(\frac{1}{t_u}\right),
\]
so, writing
\[
  L \;:=\; \log t_u, \qquad c \;:=\; \log 2\pi + 1,
\]
and absorbing all bounded quantities into a single $O(1)$,
\begin{equation}
  t_u\,(L - c) \;=\; 2\pi u + O(1). \tag{$*$}\label{eq:star}
\end{equation}

\section{The logarithmic scale of $t_u$}

\begin{lemma}\label{lem:logscale}
  As $u \to +\infty$,
  \[
    \log t_u \;=\; \log u - \log\log u + O(1).
  \]
\end{lemma}

\begin{proof}
  Both sides of \eqref{eq:star} tend to $+\infty$, and $L - c > 0$
  eventually; taking logarithms,
  \begin{equation}
    L + \log(L - c) \;=\; \log 2\pi + \log u + O\!\left(\frac1u\right).
    \tag{$**$}\label{eq:starstar}
  \end{equation}

  \emph{First pass.} Since $t_u \to \infty$ we have $L \to \infty$, and
  $\log(L-c) = o(L)$; hence \eqref{eq:starstar} forces
  $L \sim \log u$. In particular
  \[
    \log(L - c) \;=\; \log L + \log\Bigl(1 - \tfrac{c}{L}\Bigr)
    \;=\; \log L + O\!\left(\frac1L\right).
  \]

  \emph{Second pass.} From $L = (1+o(1))\log u$ we get
  \[
    \log L \;=\; \log\log u + \log(1+o(1)) \;=\; \log\log u + o(1),
  \]
  and substituting back into \eqref{eq:starstar},
  \[
    L \;=\; \log u + \log 2\pi - \log(L - c) + O\!\left(\frac1u\right)
      \;=\; \log u - \log\log u + O(1). \qedhere
  \]
\end{proof}

\section{Proof of Proposition~\ref{prop:eight} (equation (8))}

\begin{proof}
  Solving \eqref{eq:star} for $t_u$ and inserting
  Lemma~\ref{lem:logscale},
  \[
    t_u \;=\; \frac{2\pi u + O(1)}{L - c}
        \;=\; \frac{2\pi u}{\log u - \log\log u + O(1)}
              + O\!\left(\frac{1}{\log u}\right).
  \]
  Factoring out $\log u$ from the denominator and expanding
  $\frac{1}{1-x} = 1 + x + O(x^2)$ with
  $x = x_u + O\!\left(\frac{1}{\log u}\right) \to 0$,
  \[
    t_u \;=\; \frac{2\pi u}{\log u}\,
        \frac{1}{\,1 - x_u + O\!\left(\frac{1}{\log u}\right)}
        \;=\; \frac{2\pi u}{\log u}
        \left(1 + x_u + O\!\left(\frac{1}{\log u}\right)\right).
  \]
  Finally, since $\log\log u \to \infty$ we have
  $\frac{1}{\log u} = o(1)\,x_u$, so the error term
  $O\!\left(\frac{1}{\log u}\right)$ is absorbed into the coefficient of
  $x_u$:
  \[
    t_u \;=\; \frac{2\pi u}{\log u}
        \left(1 + (1 + o(1))\,\frac{\log\log u}{\log u}\right),
  \]
  which is \eqref{eq:eight}.
\end{proof}

\section{Proof of Proposition~\ref{prop:nine} (equation (9))}

\begin{proof}
  By \eqref{eq:theta2} and Lemma~\ref{lem:logscale},
  $\log\frac{t_u}{2\pi} \to \infty$, so $\vartheta'(t_u) > 0$ for all
  large $u$; since $\vartheta$ is continuously differentiable, the
  inverse function theorem shows that $u \mapsto t_u$ is differentiable,
  and differentiating \eqref{eq:seven} with respect to $u$ gives
  \[
    \vartheta'(t_u)\,\frac{dt_u}{du} \;=\; \pi,
    \qquad\text{i.e.}\qquad
    \frac{dt_u}{du} \;=\; \frac{\pi}{\vartheta'(t_u)} .
  \]

  We expand $\vartheta'(t_u)$. By \eqref{eq:theta2} and
  Lemma~\ref{lem:logscale},
  \[
    \vartheta'(t_u)
      \;=\; \frac{1}{2}\bigl(L - \log 2\pi\bigr)
            + O\!\left(\frac{1}{t_u^{2}}\right)
      \;=\; \frac{1}{2}\bigl(\log u - \log\log u + O(1)\bigr),
  \]
  where $O(t_u^{-2}) = O\!\bigl(\log^2 u / u^2\bigr)$ is far below the
  $O(1)$ already present. Inverting exactly as in the proof of
  Proposition~\ref{prop:eight},
  \[
    \frac{dt_u}{du}
      \;=\; \frac{2\pi}{\log u - \log\log u + O(1)}
      \;=\; \frac{2\pi}{\log u}
            \left(1 + x_u + O\!\left(\frac{1}{\log u}\right)\right)
      \;=\; \frac{2\pi}{\log u}
            \left(1 + (1+o(1))\,\frac{\log\log u}{\log u}\right),
  \]
  which is \eqref{eq:nine}.
\end{proof}

\begin{remark}
  The intermediate form
  $\frac{1}{\vartheta'(t_u)} = \frac{2}{\log u}
   \bigl(1 + (1+o(1))\,x_u\bigr)$
  obtained above is exactly the expansion used in the proof of
  Theorem~3 of the paper.
\end{remark}

\begin{remark}[Role of \eqref{eq:eight} and \eqref{eq:nine} in Theorem 3]
  Theorem~3 of the paper states that, for $n \ge 2$,
  \[
    t_u^{(n)} \;=\; \frac{(-1)^{n+1}\, 2\pi\,(n-2)!}{u^{n-1}\log^2 u}
      \left(1 + (2+o(1))\,\frac{\log\log u}{\log u}\right),
    \qquad u \to +\infty .
  \]
  Its proof differentiates \eqref{eq:seven} $n$ times by Fa\`a di
  Bruno's formula and isolates
  $t_u^{(n)} = -\vartheta^{(n)}(t_u)(t_u')^n / \vartheta'(t_u)$ up to
  negligible partition terms. The constant $2$ in the factor
  $(2+o(1))$ is forced by the bookkeeping of the $x_u$-coefficients of
  the three factors:
  \[
    \underbrace{t_u^{\,1-n}}_{\text{via } \eqref{eq:eight}:\; (1-n)\,x_u}
    \qquad
    \underbrace{(t_u')^{\,n}}_{\text{via } \eqref{eq:nine}:\; n\,x_u}
    \qquad
    \underbrace{1/\vartheta'(t_u)}_{\text{via the remark above}:\; x_u}
  \]
  and $(1-n) + n + 1 = 2$, independently of $n$.
\end{remark}

\end{document}
